% Section 15: Polarization
\section{Polarization in Geometric Quantization}
\label{sec:polarization}

\subsection{Introduction}

Polarization reduces the prequantum Hilbert space to the physical quantum Hilbert space by selecting half of the phase space coordinates as configuration variables.

\subsection{Kähler Polarization}

\begin{definition}[Kähler Polarization]
For the torsion field phase space with complex structure $J$, the Kähler polarization is:
\begin{equation}
    \mathcal{P} = \text{span}\left\{ \frac{\partial}{\partial a_k^*} \right\}
\end{equation}
where $a_k$ are complex coordinates defined in Section \ref{sec:prequantization}.
\end{definition}

\subsection{Polarized Sections}

Wave functions satisfying the polarization condition:
\begin{equation}
    \nabla_X \psi = 0 \quad \forall X \in \bar{\mathcal{P}}
\end{equation}
take the form:
\begin{equation}
    \psi[a, a^*] = f(a) \exp\left(-\frac{1}{2\hbar} \sum_k |a_k|^2\right)
\end{equation}

\subsection{Inner Product}

The Bargmann-Fock inner product:
\begin{equation}
    \langle \psi_1 | \psi_2 \rangle = \int \prod_k \frac{da_k \, da_k^*}{2\pi i} e^{-|a_k|^2/\hbar} \overline{f_1(a)} f_2(a)
\end{equation}

\subsection{Equivalence to Fock Space}

The map:
\begin{equation}
    f(a) \longleftrightarrow f(a^\dagger) |0\rangle
\end{equation}
establishes the unitary equivalence proved in Section \ref{sec:fock_equivalence}.
